%!TeX root=MemoriaTFG.tex

\chapter{Conclusiones}
\label{cap:conclusiones}

\section{Respuesta a la pregunta de investigación}
\label{sec:respuesta}

La pregunta formulada en sección~\ref{sec:pregunta} se respondió mediante
la evaluación empírica de catorce modelos fundacionales sobre doce
datasets externos públicos armonizados a una ontología jerárquica
L1/L2/L3, bajo cuatro protocolos principales (sondeo lineal,
ajuste fino supervisado, segmentación binaria promptable y
clasificación zero-shot multimodal) y tres protocolos
complementarios (equidad por fototipo cutáneo, comparación con
modelos de lenguaje multimodales generalistas, auditoría de
solapamiento con el dataset de preentrenamiento Derm1M). La
evidencia disponible sostiene cuatro afirmaciones cualitativas, en
las que se resumen los hallazgos defendibles del trabajo.

\paragraph{Primera. Ventaja del preentrenamiento sobre el dominio
clínico.}
PanDerm Large constituye, entre los encoders contrastados con
pesos públicos, el modelo con mejor rendimiento agregado sobre los
diez datasets externos contrastados. La mejora media respecto a
PanDerm Base es de $+4{,}1$\,pp en accuracy, $+9{,}0$\,pp en
BAcc y $+3{,}6$\,pp en AUROC bajo sondeo lineal. La mejora
absoluta es máxima sobre Dermnet ($+10{,}0$\,pp accuracy),
WSI~patches ($+8{,}7$\,pp) y PAD-UFES-20 ($+4{,}7$\,pp),
datasets con mayor número de clases o mayor heterogeneidad de
adquisición. El ajuste fino con la receta del paper original
aumenta la BAcc en $+47{,}1$\,pp sobre HAM10000 y
$+19{,}4$\,pp sobre PAD-UFES-20, atribuible al aprendizaje de las
clases minoritarias. La especialización del dominio se cuantifica
también frente a arquitecturas convolucionales supervisadas: la
ventaja de PanDerm Large sobre la mejor CNN crece con la dificultad
del dataset, desde $+0{,}5$\,pp accuracy sobre HAM10000 hasta
$+8{,}2$\,pp sobre PAD-UFES-20.

\paragraph{Segunda. Saturación temprana del sondeo lineal.}
El sondeo lineal sobre encoders preentrenados específicos de
dominio alcanza saturación con un volumen de datos etiquetados muy
reducido: con $1\%$ del conjunto de entrenamiento de HAM10000
($82$ imágenes), PanDerm Large recupera $95{,}7\%$ de la
accuracy y $96{,}2\%$ del AUROC del modelo entrenado con el
conjunto completo; con $5\%$ ($410$ imágenes), AUROC alcanza
$97{,}7\%$ del valor máximo. Sobre PAD-UFES-20, $14$ imágenes
de entrenamiento ($1\%$) bastan para alcanzar $0{,}907$
AUROC, equivalente al $95{,}5\%$ del rendimiento con el conjunto
completo. Esta eficiencia es relevante para contextos clínicos con
disponibilidad limitada de etiquetas: enfermedades raras, centros
de bajo volumen asistencial, subgrupos poco representados y
extensión a especialidades dermatológicas no cubiertas por los
benchmarks habituales.

\paragraph{Tercera. Brecha entre modelos específicos y LLMs
multimodales generalistas.}
La comparación de paradigmas sobre HAM10000 cuantifica un
gradiente de rendimiento de $43{,}5$\,pp de accuracy entre el
mejor especialista ajustado (PanDerm Base FT con TTA, $0{,}920$) y el
mejor modelo de lenguaje multimodal generalista
\emph{zero-shot} (GPT-4o, $0{,}485$). La adaptación con LoRA
sobre MedGemma~27B reduce parcialmente esta brecha
($+13{,}7$\,pp Acc HAM, $+18{,}3$\,pp Acc DDI), pero no
sistemáticamente entre datasets ($-1{,}7$\,pp sobre
PAD-UFES-20). La consecuencia operativa para sistemas de apoyo al
diagnóstico es que los LLMs multimodales generalistas no
constituyen sustitutos defendibles de los encoders fundacionales
específicos del dominio en el régimen clínico evaluado, sin
perjuicio de su papel complementario como generadores de
razonamiento textual o explicaciones.

\paragraph{Cuarta. \emph{Gap} sistemático por fototipo cutáneo.}
Los cinco encoders contrastados sobre Fitzpatrick17k completo
presentan \emph{gap}~I-VI positivo sobre la formulación de tres
clases de malignidad, con valores en BAcc entre $+0{,}124$
(BiomedCLIP) y $+0{,}306$ (DINOv2 ViT-L/14). La paradoja
BiomedCLIP, con \emph{gap} mínimo asociado a la peor BAcc global
($0{,}590$), ilustra que la métrica de equidad aislada no
constituye un criterio decisional: un modelo uniformemente inferior
puede presentar baja desigualdad sin ser equitativo en sentido
sustantivo. El modelo con mejor compromiso entre rendimiento global
y equidad es PanDerm Large (BAcc $0{,}699$, \emph{gap}
$+0{,}217$); DermLIP v2 ofrece la mejor BAcc global
($0{,}721$) con \emph{gap} superior ($+0{,}281$).

\paragraph{Síntesis. No hay modelo universalmente superior, sino
especialización por tarea.}
Leídas en conjunto, las cuatro afirmaciones anteriores no convergen
en la designación de un modelo vencedor único, sino en una lectura
matizada de la pregunta de investigación. La evidencia empírica
apunta a que el rendimiento dermatológico no es una propiedad
absoluta del codificador, sino una función conjunta del nivel
ontológico evaluado, del régimen de datos etiquetados disponibles,
de la modalidad de entrada y de la estrategia de adaptación. Cada
escenario favorece una combinación distinta de modelo y protocolo:
clasificación general, subcategoría clínica L2, cola larga L3,
ranking probabilístico, segmentación y cribado de seguridad de
melanoma no se resuelven con el mismo método óptimo. Esta lectura,
desarrollada como tesis transversal en
sección~\ref{sec:disc-especializacion}, desplaza el cuello de
botella del avance: deja de residir únicamente en la capacidad
representacional del encoder y pasa a localizarse en la taxonomía
diagnóstica, la calidad y cobertura de las etiquetas, el
alineamiento texto-imagen, la formulación de los \emph{prompts}, la
equidad por fototipo cutáneo y la estrategia de despliegue. La
respuesta a la pregunta formulada en sección~\ref{sec:pregunta} es,
por tanto, condicional: el modelo y el protocolo adecuados se
seleccionan en función del objetivo clínico y de las restricciones
de cómputo, datos y supervisión, no en abstracto.

\section{Contribuciones reproducibles}
\label{sec:contribuciones}

Las contribuciones del trabajo se enumeran en orden decreciente
de centralidad respecto a la pregunta de investigación.

\subsection{Contribuciones en el cuerpo del documento}
\label{sec:contrib-cuerpo}

\begin{enumerate}
\item \textbf{Reproducción experimental de PanDerm sobre
arquitectura aarch64.} El pipeline de PanDerm Base y Large se ha
verificado sobre NVIDIA DGX Spark con chip Grace Hopper GB10,
$128$\,GB de memoria unificada, CUDA~13.0 y precisión BF16, con
reproducción de los rangos numéricos reportados por
Yan et~al.~\cite{panderm2025} sobre los datasets compartidos. La
portabilidad al chip Grace Hopper queda documentada como base
operativa para grupos de investigación con hardware no estándar.

\item \textbf{Benchmark homogéneo de catorce modelos fundacionales
sobre doce datasets.} El estudio comparativo aplica cuatro
protocolos principales bajo configuración estándar
(sección~\ref{sec:protocolos}) sobre PanDerm Base y Large, DINOv2,
tres variantes de DermLIP, BiomedCLIP, SigLIP-Large, SAM2.1-Large,
GPT-4o, Gemini~2.5~Pro, MedGemma~4B Vision y~27B Text, y BLIP-2, más
dos líneas de base de redes convolucionales supervisadas
(ConvNeXt-Large y EfficientNetV2-Large) como referencia no
fundacional. Las cifras se reportan
con desglose por dataset y métrica
(tabla~\ref{tab:lp-panderm}, tabla~\ref{tab:benchmark},
tabla~\ref{tab:llm-comparativa}).

\item \textbf{Ontología jerárquica L1/L2/L3.} Cuatro categorías
etiológicas, $43$ subcategorías clínicas y $367$
diagnósticos individuales codificados en SNOMED~CT y CIE-10,
validados por revisión experta de la Dra.~R.~Taberner. La
ontología se aplica a la armonización de once de los doce
datasets externos contrastados, agregando $72\,654$ imágenes
sobre vocabulario común y habilitando comparaciones inter-dataset
defendibles a nivel L1 y L2.

\item \textbf{Dataset DermapixelAI~1.0.} Primer dataset
dermatológico verificado en español con imagen, metadatos y texto
narrativo asociado ($1\,089$ imágenes, $669$ casos clínicos,
mediana de $223$ palabras por caso). Particionado bajo regla
\emph{case-aware} ($891 / 157 / 41$) y distribuido bajo
licencia CC-BY-NC-SA~4.0 con datasheet formal
(Anexo~\ref{cap:anexoc}) y documento de entrega
(Anexo~\ref{cap:anexoe}).

\item \textbf{Auditoría sistemática de solapamiento con Derm1M.}
Procedimiento por hash MD5 exhaustivo sobre los doce datasets
externos contrastados, con identificación de Dermnet como único
conjunto con solapamiento exacto ($100\%$ contenido en Derm1M) y
declaración de \emph{caveat} explícito sobre las cifras zero-shot de
la familia DermLIP en ese dataset.
\end{enumerate}

\subsection{Contribuciones complementarias}
\label{sec:contrib-anexos}

\begin{enumerate}
\setcounter{enumi}{5}
\item \textbf{Análisis empírico de equidad por fototipo.}
Evaluación de cinco encoders sobre Fitzpatrick17k completo
($16\,577$ imágenes, $6$ fototipos cutáneos) con
cuantificación del \emph{gap}~I-VI por modelo
(tabla~\ref{tab:equidad-fototipo}) y análisis cruzado
\emph{Light}~$\leftrightarrow$~\emph{Dark}
(tabla~\ref{tab:cross-eval}).

\item \textbf{Caracterización de la sensibilidad al tokenizador en
DermLIP.} Identificación del tokenizador y la formulación del
\emph{prompt} como factores críticos del rendimiento zero-shot,
con diferencia de $48{,}8$\,pp AUROC entre la peor y la mejor
configuración sobre HAM10000.

\item \textbf{Adaptación de SAM2.1 al dominio dermatoscópico
mediante \emph{fine-tuning} del decodificador.} Ajuste denso del
\emph{mask\_decoder} y el \emph{promptable encoder} ($\sim 4{,}2$\,M
pesos entrenables), congelando el
\emph{image\_encoder} Hiera-Large preentrenado. Resultado de Dice
$0{,}947$ sobre ISIC2018 ($+2{,}6$\,pp sobre la cifra
publicada por Yan et~al.~\cite{panderm2025}), con transferencia
fuera de dominio prácticamente nula sobre ISIC2017 ($0{,}945$)
y PH2 ($0{,}960$).

\item \textbf{Ensemble de seguridad para detección de melanoma.}
Combinación OR \emph{top-3} de PanDerm Large ajustado sobre
HAM10000 (M1) y SigLIP-Large SO400M con sondeo lineal, que alcanza
\emph{recall} de melanoma del $100\%$
($70$ de $70$, \emph{top-3}) sobre el conjunto de test de
HAM10000.

\item \textbf{Modelo de cuello de botella conceptual SAE${}\times{}$SkinCon.}
Replicación del mecanismo de interpretabilidad por Sparse
Autoencoders sobre PanDerm Large (SAE Large,
$1\,024 \to 16\,384$ \emph{features}, entrenado sobre
$50\,454$ imágenes) y su validación frente al diccionario
clínico SkinCon~\cite{skincon} mediante un \emph{Concept
Bottleneck Model}, con AUROC media de $0{,}880$ sobre los
$22$ conceptos con cobertura suficiente y $11$ conceptos con
AUROC del mejor predictor superior a $0{,}80$. Los experimentos
E1--E4 de validación conceptual sobre Derm7pt se detallan en el
Anexo~\ref{cap:anexof}.

\item \textbf{Dermapixel R0: primer modelo fundacional
dermatológico en español.} Modelo inicializado desde el
codificador PanDerm Large y adaptado sobre el dataset
DermapixelAI~1.0, articulado en una rama supervisada (cabeza L2
con LoRA $r=16$, BAcc $0{,}363 \pm 0{,}007$ y AUROC
$0{,}855 \pm 0{,}006$, sección~\ref{sec:dermapixel-spanderm-v0})
y una rama contrastiva multimodal castellano
(sección~\ref{sec:spanderm-clip}). Constituye una versión base
entrenada sobre un corpus reducido: la contribución reside en ser
el primero de su clase y en la metodología reproducible que
establece, no en un rendimiento competitivo. Su escalado con el
banco hospitalario del HUSLL se documenta como línea de
continuación en la sección~\ref{sec:linea-spanderm}.
\end{enumerate}

\section{Hallazgos metodológicos transferibles}
\label{sec:hallazgos-conc}

Los hallazgos metodológicos del trabajo son transferibles a
escenarios clínicos análogos y constituyen aportaciones operativas
más allá del benchmark concreto evaluado. El primero es de orden
estructural ---la inexistencia de un método óptimo único--- y
enmarca a los restantes, que precisan las condiciones bajo las que
cada estrategia resulta preferible.

\paragraph{Ausencia de un método óptimo único.}
El hallazgo transversal del trabajo es que ningún codificador ni
estrategia de adaptación domina todos los escenarios: la
configuración preferible depende del nivel ontológico, del régimen
de datos, de la modalidad de entrada y del objetivo clínico
(clasificación general, subcategoría clínica L2, cola larga L3,
ranking probabilístico, segmentación o cribado de seguridad). Esta
especialización por tarea, sostenida en
sección~\ref{sec:disc-especializacion}, implica que el diseño de un
sistema de apoyo dermatológico debe formularse como selección
condicional de modelo y protocolo, y no como adopción de un único
encoder de referencia.

\paragraph{Régimen de aplicabilidad del sondeo lineal.}
La comparación cuantitativa sobre cuatro datasets
(sección~\ref{sec:disc-lp-vs-ft}) delimita un régimen operativo: para
datasets con $N_{\mathrm{train}} \lesssim 500$ imágenes, el
sondeo lineal sobre PanDerm Large supera al ajuste fino completo
en accuracy, BAcc, AUROC y F1 ponderada. El umbral se identifica
sobre DDI ($N_{\mathrm{train}} = 427$), donde el sondeo lineal
mantiene AUROC $0{,}827$ frente a $0{,}682$ del ajuste fino.

\paragraph{Compatibilidad tokenizador-encoder en zero-shot.}
La AUROC en clasificación zero-shot con modelos vision-lenguaje
contrastivos depende críticamente de la compatibilidad entre el
tokenizador empleado para los \emph{prompts} y el espacio de
\emph{embeddings} del encoder textual del modelo. La verificación
empírica de esta compatibilidad sobre un conjunto de validación
del dominio es prerrequisito de cualquier despliegue zero-shot
defendible.

\paragraph{Ajuste del decodificador sobre encoders promptables
preentrenados a gran escala.}
La combinación de SAM2.1-Large preentrenado con \emph{fine-tuning}
del decodificador de máscaras, manteniendo congelado el encoder
visual, supera al
\emph{fine-tuning} completo de baselines dermatológicos
específicos, con actualización de aproximadamente $4{,}2$\,M
parámetros entrenables (inferior al $2\%$ del total). La
generalización fuera de dominio es prácticamente nula, lo que
sugiere que las primitivas geométricas codificadas por el encoder
visual transfieren entre fuentes de adquisición sin
reentrenamiento.

\paragraph{Limitación operativa del \emph{gap} aislado como métrica
de equidad.}
La paradoja BiomedCLIP ---\emph{gap}~I-VI mínimo ($+0{,}124$ en
BAcc) asociado a la peor BAcc global ($0{,}590$)--- demuestra que la
métrica \emph{gap} aislada puede inducir conclusiones erróneas
sobre la equidad sustantiva de un modelo. El reporte conjunto de
BAcc global, AUROC global y \emph{gap} es metodológicamente
preferible y debería adoptarse como práctica estándar en estudios
de equidad dermatológica.

\paragraph{Sondeo lineal supera al ajuste fino bajo sobreajuste.}
El caso DDI ($N_{\mathrm{train}} = 427$, descenso de AUROC
$-14{,}5$\,pp con ajuste fino respecto al sondeo lineal) es
diagnóstico de la posibilidad de sobreajuste del encoder ajustado
pese a las augmentaciones de la receta canónica (Mixup, CutMix,
\emph{weighted random sampler}). La heurística operativa derivada
es que el ajuste fino sobre datasets pequeños debe ir acompañado
de monitorización de AUROC sobre conjunto de validación
independiente, y descartarse si el ranking de probabilidades
degrada.

\paragraph{Recuperación $k$-NN como alternativa competitiva en la
cola larga.}
La evidencia empírica apunta a que, sobre el nivel diagnóstico fino
L3, dominado por una distribución de cola larga estricta, la
recuperación por vecinos más próximos con índice FAISS sobre
\emph{embeddings} congelados resulta competitiva frente a las
cabezas paramétricas, que carecen de muestras suficientes por clase
para estimar fronteras de decisión estables
(sección~\ref{sec:disc-especializacion}). La heurística operativa
es que, cuando la cardinalidad de clases excede la disponibilidad de
etiquetas, la recuperación no paramétrica sobre el espacio de
representación es preferible a la clasificación supervisada directa.

\paragraph{Interpretabilidad conceptual emergente vía Sparse
Autoencoders.}
La descomposición del espacio de representación de PanDerm Large
mediante Sparse Autoencoders revela \emph{features} alineadas con
conceptos dermatológicos clínicamente interpretables sin supervisión
conceptual durante el preentrenamiento, validadas frente al
diccionario SkinCon~\cite{skincon} y frente a las anotaciones
\emph{Seven-Point Checklist} de Derm7pt. Esta interpretabilidad
emergente habilita modelos de cuello de botella conceptual que
acotan el compromiso interpretabilidad-precisión, y constituye una
vía de trazabilidad clínica transferible a otros encoders
fundacionales del dominio.

\section{Aportaciones a la práctica clínica}
\label{sec:aportaciones-clinicas}

El trabajo aporta cuatro elementos directamente aplicables al
diseño de sistemas de apoyo al diagnóstico dermatológico en el
contexto del Servicio de Dermatología del HUSLL y, por extensión,
al circuito de teledermatología institucional descrito en
sección~\ref{sec:trabajo-previo}.

\paragraph{Arquitectura técnica defendible.}
Encoder PanDerm Large preentrenado y congelado por defecto, con
cabeza lineal multi-clase mediante regresión logística L-BFGS
($C = 1{,}0$, \texttt{max\_iter}~$=5\,000$) en datasets con
$N_{\mathrm{train}} \lesssim 500$, o cabeza ajustada con la
receta de PanDerm Large en datasets con
$N_{\mathrm{train}} \gtrsim 5\,000$ y desbalance acusado.
Segmentación con SAM2.1-Large adaptado mediante \emph{fine-tuning}
del decodificador cuando la tarea requiera localización de lesión. Ensemble OR \emph{top-3}
de M1 (PanDerm Large \emph{fine-tuned}) y SigLIP-Large con sondeo
lineal cuando la tarea sea cribado oportunista de melanoma con
prioridad de \emph{recall}.

\paragraph{Ontología armonizada como columna vertebral.}
La ontología L1/L2/L3 con codificación SNOMED~CT y CIE-10
proporciona una columna vertebral semántica reutilizable para
sistemas de información hospitalarios y compatible con la
nomenclatura clínica del entorno. La presentación jerárquica al
usuario clínico (predicción L3 acompañada de su agregación L2 y
L1) aporta robustez frente a la incertidumbre en la cola larga sin
sacrificar especificidad cuando ésta sea fiable.

\paragraph{Dataset DermapixelAI~1.0 como recurso público en español.}
Habilita la reutilización del archivo Dermapixel en proyectos
académicos posteriores, la adaptación de modelos vision-lenguaje a
texto clínico en castellano y la validación de sistemas de apoyo
sobre material de referencia docente reconocida en el ámbito
hispanohablante.

\paragraph{Caveats explícitos sobre fuga de información.}
La auditoría de solapamiento entre datasets de evaluación y
preentrenamiento aporta una práctica de reporte transferible a
otros benchmarks dermatológicos. La declaración explícita de
\emph{caveats} de leakage en las tablas
(tabla~\ref{tab:datasets-resumen}) y la consistencia entre la cifra
reportada y su \emph{caveat} asociado debería adoptarse como
condición de defensibilidad en publicaciones del campo.

\section{Líneas de investigación abiertas}
\label{sec:lineas-abiertas}

Las líneas abiertas que se desarrollan a continuación constituyen
propuestas concretas con experimentos identificados, prerrequisitos
declarados y métricas esperables. No son trabajo simplemente
pendiente, sino aportaciones intelectuales del proyecto cuya
ejecución requiere recursos o aprobaciones específicas. Todas
descansan sobre las contribuciones reproducibles enumeradas en
sección~\ref{sec:contribuciones}.

\subsection{Cobertura experimental sobre DermapixelAI~1.0: técnicas ejecutadas y pendientes}
\label{sec:dermapixel-futuro}

Las secciones~\ref{sec:dermapixel-results}
y~\ref{sec:dermapixel-zs} del capítulo~\ref{cap:resultados} evalúan
PanDerm Base, PanDerm Large, DermLIP v2, SigLIP-SO400M y BiomedCLIP
sobre DermapixelAI~1.0 exclusivamente como extractores de
\emph{features} congelados, en combinación con cabezas paramétricas
ligeras (regresión logística, $k$-NN, MLP de $512$ unidades) o con
clasificación contrastiva \emph{zero-shot}. El espacio de hipótesis
no explorado sobre el dataset en castellano incluye un conjunto de
técnicas de \emph{transfer learning} denso que la literatura sobre
modelos fundacionales en imagen médica documenta como aportando
incrementos típicos de $10$ a $20$\,pp en BAcc respecto al sondeo
lineal congelado, en línea con los $+47{,}1$\,pp BAcc sobre
HAM10000 y $+19{,}4$\,pp BAcc sobre PAD-UFES-20 cuantificados en
la sección~\ref{sec:ft-results}. La presente sección documenta de
forma explícita estas técnicas pendientes y la cifra esperable
sobre L2 (38 clases) tras su aplicación, condicionada a la
disponibilidad del banco hospitalario HUSLL como ampliación del
dataset.

\begin{table}[H]
\centering
\footnotesize
\setlength{\tabcolsep}{4pt}
\renewcommand{\arraystretch}{0.95}
\caption{Técnicas de \emph{transfer learning} pendientes de
aplicación sobre DermapixelAI~1.0 con el cuerpo experimental del
presente trabajo. \emph{Coste}: orden de magnitud sobre NVIDIA DGX
Spark GB10. \emph{Esperado}: mejora cualitativa estimada a partir
de las referencias internas del trabajo. \emph{Estado}: situación
respecto al cierre del presente TFG.}
\label{tab:dermapixel-pendiente}
\begin{tabular}{p{0.19\textwidth}p{0.09\textwidth}p{0.15\textwidth}p{0.30\textwidth}p{0.15\textwidth}}
\hline
Método & Tipo & Coste estimado & Mejora esperada & Estado \\
\hline
Ajuste fino completo del encoder        & Supervisado            & $\sim 5$ min Spark (3 niveles) & $+18{,}2$\,pp BAcc L1 sobre LP CV; $-4{,}0$\,pp BAcc L2 frente a Dermapixel R0 LoRA; sin mejora sobre LP en L3 & \textbf{Ejecutado} (L1$+$L2$+$L3), sección~\ref{sec:dermapixel-ft} \\
LoRA $r=16$ (Dermapixel R0)             & PEFT                   & $\sim 12$ min Spark (3 seeds) & $+19{,}5$\,pp BAcc L2 sobre LP CV; AUROC L2 $0{,}855$ empata con DermLIP individual & \textbf{Ejecutado}, sección~\ref{sec:dermapixel-spanderm-v0} \\
\emph{Test-time augmentation}           & Inferencia             & $5$--$10$ min              & $+3$ a $+22$\,pp Acc@3 (cf.\ sección~\ref{sec:ensemble-results}) & \textbf{Probado}, sección~\ref{sec:dermapixel-tta} \\
Ensemble Base$+$Large$+$DermLIP         & Inferencia             & Trivial                    & $+6{,}7$\,pp AUROC L2 con DermLIP individual; $+2{,}1$\,pp AUROC L1 con avg Large$+$DermLIP & \textbf{Probado}, sección~\ref{sec:dermapixel-ensemble} \\
Pérdida focal / re-ponderación de clases & Entrenamiento         & $\sim 5$ min Spark         & $+3{,}3$\,pp BAcc L1 y $+5{,}9$\,pp BAcc L2 sobre LP CE ($\gamma = 2{,}0$); sin efecto L3 & \textbf{Ejecutado} ($\gamma \in \{1,\,2,\,5\} \times \{\text{none},\,\text{balanced}\}$), sección~\ref{sec:dermapixel-focal} \\
CBM sobre conceptos SkinCon              & Interpretable          & Anotación dermatológica    & Trazabilidad clínica                                          & Diseñado (Anexo~\ref{cap:anexof}) \\
Banco hospitalario HUSLL ($\sim 22$\,K) & Datos                  & $3$--$6$ meses (CEIC)      & Activación L3 directa (cabeza L3 de Dermapixel R0)            & Bloqueado (revisión ética) \\
\hline
\end{tabular}
\end{table}

La elección de evaluar sobre encoder congelado en las
secciones~\ref{sec:dermapixel-results} y~\ref{sec:dermapixel-zs}
sigue el protocolo estándar de la literatura sobre modelos
fundacionales en imagen médica, que prioriza la cuantificación
limpia del valor del preentrenamiento sobre datos del dominio sin
contaminación por hiperparámetros específicos de ajuste fino. La
familia PanDerm~\cite{panderm2025}, DermLIP~\cite{derm1m} y
BiomedCLIP~\cite{biomedclip} reportan sus cifras principales bajo
este protocolo, lo que permite la comparación directa con la
literatura. Los métodos pendientes recogidos en la
tabla~\ref{tab:dermapixel-pendiente} se concentran tras el cierre
experimental del presente trabajo en la extensión del dataset vía
banco hospitalario HUSLL (bloqueada por la revisión ética del
CEIC) y en la línea de modelos de cuello de botella conceptual
(CBM) sobre conceptos SkinCon, condicionada a la anotación
dermatológica detallada en el Anexo~\ref{cap:anexof}. Las
restantes técnicas (ajuste fino denso L1$+$L2$+$L3,
adaptación LoRA de Dermapixel R0, \emph{Test-Time Augmentation},
ensemble multi-codificador y pérdida focal) han sido ejecutadas y
reportadas en las secciones correspondientes del
capítulo~\ref{cap:resultados}.

Bajo la hipótesis de aplicación combinada de las técnicas
pendientes (ajuste fino del encoder con la receta de la
sección~\ref{sec:ft-results}, ensemble PanDerm$+$DermLIP en la
línea documentada en la sección~\ref{sec:ensemble-results} y
ampliación con el banco hospitalario HUSLL hasta superar el umbral
de aproximadamente $30$ imágenes por clase L3 que activa la
variante con cabeza L3 directa de Dermapixel R0), las cifras esperables para un
clasificador desplegable sobre L2 ($38$ subcategorías clínicas)
en castellano se estiman en BAcc entre $0{,}45$ y $0{,}55$,
orden de magnitud comparable con el rango reportado en la
sección~\ref{sec:ft-results} para PanDerm Base ajustado sobre
HAM10000 (BAcc $0{,}575$ sobre $7$ clases). Esta estimación es
prospectiva, no constituye un compromiso ni una predicción puntual,
y queda condicionada a la disponibilidad efectiva del banco
hospitalario tras la revisión por el comité ético de investigación
clínica y a la calibración empírica de la receta de ajuste fino
sobre el material en castellano.

\subsection{Dermapixel R0: primer modelo fundacional dermatológico
en español}
\label{sec:linea-spanderm}

El dataset DermapixelAI~1.0 habilita el desarrollo de un modelo
fundacional dermatológico clínico en español, denominado
\textbf{Dermapixel R0}, inicializado desde el codificador PanDerm
Large y adaptado sobre dicho dataset. La denominación R0 designa
una versión base entrenada con un corpus reducido: un punto de
partida que sienta las bases para iteraciones posteriores, no un
modelo terminado. El modelo se articula en dos ramas
complementarias.

La primera es la rama \textbf{supervisada}: cabeza clasificadora
sobre encoder PanDerm Large adaptado mediante \emph{Low-Rank
Adaptation}. Comprende la cabeza L2 con LoRA $r=16$
ejecutada y reportada en la
sección~\ref{sec:dermapixel-spanderm-v0} y una cabeza L3 directa
---condicionada al banco hospitalario HUSLL---,
descritas a continuación en los párrafos correspondientes.

La segunda es la rama \textbf{zero-shot multimodal}:
la rama contrastiva de Dermapixel R0, sistema dual-encoder estilo CLIP que alinea el
encoder visual PanDerm Large con un encoder textual multilingüe
(\texttt{intfloat/multilingual-e5-base}) sobre el mismo dataset
DermapixelAI~1.0. Se han ejecutado cinco iteraciones
experimentales (v1 a v5) reportadas íntegramente en la
sección~\ref{sec:spanderm-clip}. La iteración v4
($\text{LoRA}$ visual $r=32$ sobre los últimos doce bloques del
ViT, régimen híbrido \emph{H6} en inferencia) alcanza L3
\emph{accuracy}@1 \emph{zero-shot} $= 0{,}2475$ sobre el conjunto
de test de $101$ imágenes, esencialmente el umbral paper-grade
$0{,}25$ fijado a priori sobre vocabulario castellano fino. La
iteración v5, bajo configuración LoRA idéntica y ampliación del
dataset con $15\,128$ imágenes externas genéricas, reporta una
ablación negativa documentada ($0{,}2079$, $-4{,}0$\,pp): la
escala genérica mejora las representaciones visuales aisladas
pero erosiona la alineación cross-modal sobre el registro
discursivo castellano. La rama aporta tres resultados empíricos
complementarios. El primero es la recuperación cross-modal
castellano-imagen, con un \emph{lift} de $\mathbf{14{,}0\times}$
sobre el azar (v2, $i2t$ R@$5 = 0{,}654$ frente a $0{,}047$ en
régimen de \emph{caption} rica, \emph{pool} $= 107$). El segundo
es el diagnóstico mecanístico del techo del encoder visual
congelado mediante \emph{Linear Probing}. El tercero es la
reducción del ECE L3 en un $97\%$
($0{,}7262 \to 0{,}0213$ con $\tau = 0{,}4165$ sobre v4) mediante
\emph{temperature scaling}~\cite{guo2017}. Una ablación
independiente con un modelo fundacional inglés grande preentrenado
confirma que el cuello de botella experimental era la escala de la
alineación cross-modal castellano, no el codificador visual ni la
formulación textual (Anexo~\ref{cap:anexoi}). La rama zero-shot no es sustitutiva de
la supervisada: ambas operan sobre regímenes de uso distintos
(vocabulario abierto frente a vocabulario cerrado L2/L3) y la
integración del prototipo DermApIxel (Anexo~\ref{cap:anexoh})
las expone como módulos paralelos del pipeline.

Los párrafos siguientes describen los prerrequisitos y el estado
de las dos versiones de la rama supervisada.

\paragraph{Dermapixel R0 (cabeza supervisada L2).}
Cabeza L2 ($38$ clases efectivas tras consolidación de
variantes ortográficas) sobre encoder PanDerm Large adaptado
mediante \emph{Low-Rank Adaptation} (LoRA, $r=16$, $\alpha=32$)
sobre los dos últimos bloques de transformer, con ranking L3 por
similitud coseno sobre los \emph{embeddings} prototípicos del
conjunto de entrenamiento de DermapixelAI~1.0. El protocolo es
ejecutable sobre el hardware actual del proyecto sin recursos
adicionales. La métrica primaria prevista era Acc-top-1 sobre L2
y \emph{recall}-top-3 sobre L3 en el conjunto de test, con
validación cruzada estratificada por L1 dada la limitación del
tamaño muestral.

La ejecución del protocolo (sección~\ref{sec:dermapixel-spanderm-v0}
del capítulo~\ref{cap:resultados}) reporta sobre el nivel L2
BAcc $= 0{,}363 \pm 0{,}007$, AUROC $= 0{,}855 \pm 0{,}006$ y
Acc@3 $= 0{,}500 \pm 0{,}023$ (media $\pm$ desviación estándar
sobre tres semillas, con $524\,K$ parámetros entrenables que
representan el $0{,}17\%$ del encoder). Sobre el ranking L3 por
prototipos, Acc@1 $= 0{,}167$, Acc@3 $= 0{,}259 \pm 0{,}013$ y
Acc@5 $= 0{,}296 \pm 0{,}026$ sobre $250$ clases. Las cifras
observadas alcanzan o superan las expectativas del diseño: la
BAcc L2 mejora en $+19{,}5$\,pp la cifra de validación cruzada
del sondeo lineal congelado y la AUROC L2 empata
estadísticamente con la de DermLIP v2 individual, lo que valida
empíricamente la línea LoRA como protocolo de referencia sobre el
dataset en castellano.

\paragraph{Dermapixel R0 (cabeza L3 directa).}
Cabeza L3 directa ($251$ diagnósticos efectivos sobre el
vocabulario completo) condicionada a la incorporación del banco
hospitalario del HUSLL ($\sim 22\,000$ estudios estimados tras
aprobación CEIC). El umbral de activación es de al menos
$30$ imágenes por clase L3, condición no satisfecha por
DermapixelAI~1.0 en aislamiento dado que $178$ entradas L3 del
vocabulario con muestras efectivas contienen $\leq 5$ imágenes
(cola larga estricta del vocabulario representado). La extensión vía
banco HUSLL es condición necesaria y suficiente para la viabilidad
de esta cabeza L3 directa.

\subsection{Derm7pt + Sparse Autoencoders: ground-truth conceptual
para interpretabilidad}
\label{sec:linea-derm7pt-sae}

Las anotaciones \emph{Seven-Point Checklist} de Derm7pt
(siete criterios diagnósticos estructurados por imagen) constituyen
ground-truth conceptual no explotada por la literatura previa sobre
Sparse Autoencoders aplicados a dermatología. La propuesta de
investigación contempla cuatro experimentos. Los tres primeros se
han ejecutado y se reportan en la
sección~\ref{sec:sae-derm7pt}; el cuarto se ha ejecutado y se
reporta en la sección~\ref{sec:sae-d7pt-e4} dentro de la misma
sección de resultados.

\paragraph{E1. Correlación SAE-Derm7pt. \emph{Ejecutado}.}
Cálculo de la AUROC efectiva de cada \emph{feature} del SAE Large
($16\,384$ \emph{features}) como predictor binario de cada uno
de los siete criterios Derm7pt sobre las anotaciones de la variante
dermatoscópica ($N = 1\,011$). El resultado, reportado en
la tabla~\ref{tab:sae-d7pt-e1} de la
sección~\ref{sec:sae-derm7pt}, identifica AUROC efectivo
\emph{top-1} hasta $0{,}823$ sobre \emph{vascular structures}
y entre $0{,}681$ y $0{,}754$ sobre los seis criterios
restantes.

\paragraph{E2. \emph{Concept Bottleneck Model} sobre Derm7pt.
\emph{Ejecutado}.}
Construcción de un CBM jerárquico análogo al utilizado sobre
SkinCon~\cite{skincon}, con los siete criterios Derm7pt como
cuello de botella conceptual entre encoder y clasificador binario
melanoma/no-melanoma sobre los \emph{splits} oficiales
(\emph{train}~$=413$, \emph{val}~$=203$, \emph{test}~$=395$). La
comparativa con el sondeo lineal directo
(tabla~\ref{tab:sae-d7pt-cbm} de la
sección~\ref{sec:sae-derm7pt}) cuantifica el compromiso
interpretabilidad-precisión en $-5{,}8$\,pp AUROC
($0{,}890 \to 0{,}832$) sobre el conjunto de test. Los pesos del
meta-clasificador (tabla~\ref{tab:sae-d7pt-cbm-coefs})
priorizan \emph{blue whitish veil} y \emph{regression structures}
de forma coherente con la escala clínica.

\paragraph{E3. Cruce con conceptos del Grupo~A. \emph{Ejecutado}.}
Mapeo operativo de los $16$ conceptos dermatoscópicos
propuestos a la Dra.~R.~Taberner sobre los siete criterios del
\emph{Seven-Point Checklist} y validación de la cobertura del
SAE Large sobre cada concepto mapeado. La
tabla~\ref{tab:sae-d7pt-cross-rosa} de la
sección~\ref{sec:sae-derm7pt} documenta que $11$ de
$16$ conceptos admiten \emph{mapping} y $7$ de los $11$
quedan cubiertos por al menos una \emph{feature} SAE con AUROC
$\geq 0{,}70$. El pico observado sobre vasos en horquilla
($0{,}926$, $N_{+} = 15$) sostiene la línea de
interpretabilidad sin entrenamiento específico sobre el dataset en
castellano. Los $5$ conceptos sin \emph{mapping} (empedrado,
paralelos surcos/crestas, lagunas vasculares, vasos en corona) son
estructuras no contempladas en la escala original y constituyen
precisamente el contenido específico de la anotación clínica
colaborativa con la Dra.~R.~Taberner.

\paragraph{E4. \emph{Fine-tuning} multitarea \emph{concept-aware}
sobre Derm7pt. \emph{Ejecutado}.}
Ajuste fino supervisado de PanDerm Large con adaptación de bajo
rango (LoRA, $r = 16$) sobre los dos últimos bloques del
codificador y ocho cabezas paralelas: una multi-clase por cada
criterio del \emph{Seven-Point Checklist} y una binaria para
melanoma/no-melanoma. La pérdida es la suma equiponderada de
entropías cruzadas sobre las ocho cabezas. La configuración,
los hiperparámetros y los resultados se reportan en la
sección~\ref{sec:sae-d7pt-e4}: con sólo $0{,}56$\,M de
parámetros entrenables ($0{,}18\%$ del total) la cabeza
melanoma alcanza AUROC test $= 0{,}891$ (IC95\%
$[0{,}856,\,0{,}927]$), superando marginalmente al sondeo lineal
directo sobre el SAE ($0{,}890$) y mejorando el CBM jerárquico
en $+5{,}9$\,pp AUROC. Las siete cabezas conceptuales pierden
entre $1$ y $3$\,pp de AUROC respecto al sondeo lineal separado
por concepto, coste compatible con el comportamiento típico del
aprendizaje multitarea y aceptado a cambio de inferencia
simultánea melanoma + siete criterios en una sola pasada. Este
resultado consolida la integración del módulo M7 de
interpretabilidad del prototipo DermApIxel
(Anexo~\ref{cap:anexoh}) sobre una única arquitectura adaptada y
elimina la dependencia de pipelines en cascada para producir
inferencia conceptual.

\subsection{Aprendizaje activo \emph{in-domain} sobre clases L3
con peor cobertura}
\label{sec:linea-active-learning}

La trilogía experimental de la rama contrastiva de Dermapixel R0
reportada en
sección~\ref{sec:spanderm-clip} delimita con precisión la palanca
correcta para superar el umbral L3 \emph{accuracy}@1
$\sim 0{,}25$ bajo \emph{zero-shot} castellano. La iteración v4
($\text{LoRA}$ visual $r=32$ sobre los últimos doce bloques del
ViT, régimen híbrido \emph{H6} en inferencia) alcanza
$0{,}2475$ sobre el conjunto de test fijo, esencialmente el
umbral fijado a priori. La iteración v5, bajo configuración
arquitectónica idéntica y ampliación del dataset de
entrenamiento con $15\,128$ imágenes externas genéricas, queda
\emph{flat-to-negative} ($0{,}2079$): la escala genérica
mejora las representaciones visuales aisladas ($+41{,}2\%$
relativo sobre régimen visual puro \emph{H1}, $+54{,}4\%$
relativo sobre techo \emph{Linear Probing} L2 supervisado) pero
erosiona la alineación cross-modal sobre el registro discursivo
castellano del archivo original. La cola larga estructural del nivel L3
permanece sin cobertura: $327$ clases con $\leq 5$ imágenes sobre
el vocabulario ontológico completo ($178$ con muestras $1$--$5$ en
entrenamiento, $149$ sin imágenes en el dataset actual) no se
rellenan tras la incorporación del dataset externo.

La línea identificada como palanca correcta para superar
$0{,}25$ bajo \emph{zero-shot} castellano fino consiste en una
estrategia de aprendizaje activo dirigida \emph{in-domain}: en
lugar de ampliar el dataset de entrenamiento de forma
indiscriminada, se identifican las clases L3 con peor cobertura
en el subconjunto castellano (las ocho clases críticas
$\emph{deepened}$ en v5 ---seis tumorales más dermatitis atópica
y urticaria--- como muestra inicial extensible) y se invierte
esfuerzo experto en anotación dirigida sobre estas clases. Los
candidatos a anotación se generan mediante consulta del
prototipo DermApIxel sobre cohortes hospitalarias del HUSLL bajo
revisión por la Dra.~R.~Taberner. El banco hospitalario HUSLL
($\sim 22\,000$ estudios estimados tras aprobación CEIC,
sección~\ref{sec:linea-husll}) constituye el recurso natural
para esta estrategia: su composición clínica refleja la
distribución diagnóstica real del Servicio de Dermatología, su
material está alineado con el registro castellano del archivo
Dermapixel~1.0 y su escala permite tanto el llenado dirigido de
la cola larga L3 como la eventual activación del modelo
supervisado directo L3 (cabeza L3 de Dermapixel R0) cuando se supere el umbral
operativo de $\sim 30$ imágenes por clase L3
(sección~\ref{sec:linea-spanderm}). La consecución del umbral
operativo $0{,}25$ por v4 sobre el dataset actual hace que esta
línea de continuación ya no se plantee como ruta hacia el
\emph{target}, sino como ruta de generalización del sistema
hacia cobertura clínica completa. Esta dirección queda reforzada
por la ablación del Anexo~\ref{cap:anexoi}: la palanca pendiente
para la rama contrastiva de Dermapixel R0 no es escala genérica de datos sino anotación
experta \emph{in-domain} sobre las clases L3 con peor cobertura.

\subsection{Selección de arquitectura para segmentación promptable}
\label{sec:linea-seg-arquitectura}

La comparativa pareada de arquitecturas promptables
(sección~\ref{sec:seg-arquitecturas}) extiende el resultado de
segmentación de la sección~\ref{sec:seg-results} sobre el mismo
test canónico de ISIC2018 y delimita una pauta de selección. La
generación de arquitectura SAM2 domina sobre SAM v1 con
independencia del preentrenamiento de dominio, el preentrenamiento
médico aporta una mejora consistente solo dentro de la misma
generación, y el tamaño de \emph{backbone} no aporta ventaja en el
régimen \emph{box-prompted}. La configuración con mejor Dice
($0{,}9556$) y latencia próxima a la mínima ($26$\,ms) es
MedSAM2-tiny, que se identifica como candidata de referencia para
la segmentación de lesión en despliegues con restricciones de
cómputo. Esta selección se acompaña de tres comprobaciones que
consolidan la recomendación
(secciones~\ref{sec:seg-control-automatico}--\ref{sec:seg-techo-anotador}):
el factor dominante del rendimiento es la caja-\emph{prompt} y no el
segmentador ---retirarla cuesta $7{,}5$\,pp frente a los $\sim 1$\,pp
que separan a las arquitecturas---; la cifra es estable bajo
validación cruzada de cinco pliegues
($0{,}9554\pm 0{,}0010$) y la adaptación de bajo rango
(LoRA~\cite{lora}) iguala al ajuste denso sin ventaja; y el resultado
se sitúa cerca del techo de acuerdo inter-anotador ($0{,}728$ de Dice
entre humanos frente a $0{,}915$ del modelo contra el consenso). La
asimetría de la robustez a la caja ---apretarla degrada mucho más que
aflojarla--- fija además la especificación del detector para el modo
automático. Quedan abiertas la evaluación con detector de lesión
extremo a extremo y la validación prospectiva sobre adquisición
clínica no dermatoscópica.

\subsection{Escalado del dataset con el banco hospitalario HUSLL}
\label{sec:linea-husll}

La extensión del dataset DermapixelAI~1.0 con el banco hospitalario
del HUSLL constituye la línea de mayor impacto potencial sobre el
rendimiento clínico real del sistema. La estimación preliminar es
de aproximadamente $22\,000$ estudios disponibles en archivo,
con metadatos asociados (imagen clínica, dermatoscópica cuando
disponible, diagnóstico final, edad y sexo del paciente). Los
prerrequisitos son cuatro: aprobación del Comité Ético de
Investigación Clínica de las Illes Balears (CEIC-IB),
implementación de un protocolo de anonimización compatible con
RGPD y la Ley Orgánica $3/2018$, definición de un esquema de
licencia compatible con la naturaleza hospitalaria del material
(probablemente más restrictiva que CC-BY-NC-SA~4.0), y
establecimiento de un equipo de validación experta liderado por la
Dra.~R.~Taberner con apoyo de residentes de dermatología del HUSLL.

La incorporación del banco hospitalario habilita la cabeza L3 directa de Dermapixel R0
(sección~\ref{sec:linea-spanderm}), la validación clínica prospectiva
sobre cohorte real (sección~\ref{sec:linea-prospectivo}) y un análisis
de equidad sobre fototipo cutáneo más representativo de la
población atendida por el sistema sanitario español que el
disponible en Fitzpatrick17k.

\subsection{Validación clínica prospectiva}
\label{sec:linea-prospectivo}

La validación clínica prospectiva del prototipo DermApIxel
(Anexo~\ref{cap:anexoh}) sobre cohorte hospitalaria real
del HUSLL requiere un protocolo aprobado por CEIC-IB con cuatro
elementos: definición de \emph{endpoints} clínicos primarios
(\emph{recall} de melanoma, tasa de derivaciones evitadas) y
secundarios (concordancia inter-observador, tiempo de
diagnóstico, valoración subjetiva del clínico); cálculo
muestral basado en los \emph{recall} preliminares de
$\sim 0{,}97$ obtenidos sobre HAM10000 con M1+SigLIP y un
\emph{estándar de oro} conservador; reclutamiento de dermatólogos
revisores independientes para arbitraje en casos de discrepancia;
y plan de análisis estadístico preregistrado.

La cifra de \emph{recall} del $100\%$ sobre melanoma reportada
en sección~\ref{sec:ensemble-results} es prueba de concepto sobre
material de archivo, no resultado clínico validado. Su
extrapolación al flujo asistencial real requiere la validación
prospectiva aquí esbozada, identificada como prerrequisito de
cualquier despliegue operativo más allá del prototipo.

\subsection{Auditoría sistemática de leakage extensible}
\label{sec:linea-leakage}

La auditoría por hash MD5 desarrollada para Derm1M
(sección~\ref{sec:auditoria}) es directamente extensible a los demás
corpora de preentrenamiento del ecosistema dermatológico actual:
SkinGPT-4, MONET, MM-Skin, MedSAM y SkinFM, entre otros. La
propuesta consiste en aplicar el mismo procedimiento sobre los
splits de evaluación canónicos contrastados en este trabajo y
publicar las cardinalidades de intersección como recurso de
referencia para la comunidad. El producto previsto es un reporte
técnico con la tabla completa de solapamientos por (dataset de
evaluación, dataset de preentrenamiento), acompañado del script
auditable. El \emph{caveat} cuantificado por dataset permitiría
recalibrar las cifras zero-shot publicadas sobre estos corpora.

\subsection{Recuperación visual densa como sistema de apoyo}
\label{sec:linea-rag}

La recuperación visual densa basada en FAISS sobre los
\emph{embeddings} de Derm1M ($\sim 421\,000$ imágenes
preindexadas) constituye un componente del prototipo DermApIxel
(sección~\ref{sec:dermapixel-m4} del Anexo~\ref{cap:anexoh}) con
latencia inferior a $150$
milisegundos por consulta. La propuesta consiste en evaluar el
valor clínico del módulo mediante un estudio mixto sobre cohorte
HUSLL: cuantitativo (concordancia diagnóstico-recuperación,
\emph{precision@k} sobre L1 y L2 de la ontología) y cualitativo
(valoración subjetiva del clínico sobre la utilidad de los casos
recuperados como referencia diferencial). La integración con la
historia clínica electrónica del HUSLL queda como prerrequisito
técnico no abordado en el presente trabajo.

\subsection{IA agéntica para razonamiento clínico estructurado}
\label{sec:linea-agentica}

A medio plazo, la combinación de los componentes desarrollados
(encoder PanDerm, ontología L1/L2/L3, dataset DermapixelAI,
recuperación visual densa y LLM multimodal en local) habilita la
construcción de un agente clínico orquestado bajo control del
facultativo, que incorpore evidencia estructurada externa y
auditable. Frente al uso directo de un LLM multimodal generalista
(sección~\ref{sec:disc-llm}), este esquema reduciría el sesgo
léxico documentado y facilitaría la justificación clínica de la
salida. Su desarrollo es condicional a la disponibilidad de pesos
abiertos con capacidad de razonamiento estructurado y se detalla en
el Anexo~\ref{cap:anexoh}.

\subsection{Reentrenamiento multilingüe de la rama lingüística}
\label{sec:linea-multilingual}

La sensibilidad documentada al tokenizador
(sección~\ref{sec:disc-zs}) sugiere que un reentrenamiento dirigido del
encoder textual de DermLIP sobre texto clínico en español aportaría
mejoras apreciables sobre el rendimiento zero-shot en el ámbito
hispanohablante. La estrategia mínima viable consiste en aplicar
\emph{continued pretraining} del encoder textual sobre los
$669$ casos clínicos de DermapixelAI~1.0 con el razonamiento
narrativo asociado, manteniendo congelado el encoder visual de
DermLIP. La estrategia ampliada requiere un dataset de pares
imagen-texto en español de mayor escala, no disponible en la fecha
de cierre del trabajo y condicional al escalado descrito en
sección~\ref{sec:linea-husll}.

\section{Cierre}
\label{sec:cierre}

El trabajo ha caracterizado empíricamente el rendimiento y la
capacidad de generalización de catorce modelos fundacionales para
imagen dermatológica disponibles, bajo cuatro protocolos
principales y tres complementarios, sobre doce datasets externos
armonizados a una ontología jerárquica de tres niveles validada
por revisión experta. La evidencia obtenida sostiene cuatro
afirmaciones cualitativas defendibles (sección~\ref{sec:respuesta}),
las contribuciones reproducibles enumeradas en
sección~\ref{sec:contribuciones}, los hallazgos metodológicos
transferibles de la sección~\ref{sec:hallazgos-conc} y cuatro
aportaciones a la práctica clínica
(sección~\ref{sec:aportaciones-clinicas}).

La lectura de conjunto de esta evidencia no es la designación de un
codificador vencedor, sino una tesis matizada: el rendimiento
dermatológico depende del nivel ontológico, del régimen de datos,
de la modalidad de entrada y de la estrategia de adaptación, de modo
que la elección de modelo y protocolo se resuelve en función del
objetivo clínico y no en abstracto
(sección~\ref{sec:disc-especializacion}). Esta lectura matiza
cualquier ordenación global de los modelos evaluados y desplaza la
palanca de mejora desde la capacidad representacional del encoder
hacia la taxonomía diagnóstica, la cobertura de etiquetas, el
alineamiento texto-imagen, los \emph{prompts}, la equidad por
fototipo cutáneo y la estrategia de despliegue.

Las líneas de investigación abiertas
(sección~\ref{sec:lineas-abiertas}) constituyen propuestas concretas con
experimentos identificados y prerrequisitos declarados, no trabajo
simplemente pendiente. Su ejecución completaría la cadena entre
caracterización empírica del estado del arte, construcción de
recursos públicos en español, validación clínica prospectiva y
despliegue operativo en el sistema sanitario, recorrido cuya
primera etapa cierra el presente documento.

La aportación intelectual del trabajo se sitúa en la conjunción de
tres elementos. El primero es una contribución empírica: la
evaluación rigurosa y multi-eje del estado del arte, que cuantifica
la dependencia del rendimiento respecto al nivel ontológico, el
régimen de datos, la modalidad y la estrategia de adaptación. El
segundo es una contribución de recurso: el dataset DermapixelAI~1.0
y la semilla Dermapixel R0
(contribuciones~\ref{sec:contrib-cuerpo}
y~\ref{sec:contrib-anexos}), primeros de su clase para texto
clínico dermatológico en español. El tercero es la formulación
explícita de líneas de continuación con base técnica y clínica
verificable. De esta evidencia se desprende que el avance del campo
no pasa por un único encoder de referencia, sino por una
arquitectura modular que combine codificadores, niveles ontológicos
y estrategias de adaptación según la tarea, junto con decisiones de
despliegue que contemplen cobertura de etiquetas, equidad y
trazabilidad. El siguiente paso natural es la primera publicación
derivada del trabajo, centrada en la comparativa de modelos
fundacionales bajo la ontología armonizada, seguida del escalado de
Dermapixel R0 como modelo fundacional dermatológico para texto
clínico en español a partir del dataset DermapixelAI~1.0 extendido
con el banco hospitalario del HUSLL.
